% GENERATED FILE. Edit canonical lessons, then run npm run generate:books.
% canonical-source-hash: fnv1a64:c40be778ec8453a7
% canonical-lessons: PA-C278-hisab, PA-C278-laddna, PA-C278-vikna, PA-C278-paikat, PA-C278-jorna

\chapter{Packets and Accounts}
\label{ch:pa-a2-packets-and-accounts}
\hlchaptermodality{278}

\begin{chapteropening}
\textbf{By the end of this chapter:} \emph{I can talk about goods that are packed and sold.}

Complete the last lesson of chapter 278: I can talk about goods that are packed and sold.
\end{chapteropening}

\section[\texorpdfstring{\pa{ਹਿਸਾਬ} (hisāb)}{ਹਿਸਾਬ (hisāb)}]{\pa{ਹਿਸਾਬ} (hisāb) — a reckoning; an account of money}

\label{lesson:PA-C278-hisab}



\begin{quote}
Before the new one: say the Punjabi for a thing, then the Punjabi for goods.
\end{quote}

\subsection*{You'll want to know: \pa{ਹਿਸਾਬ}}
\textbf{\pa{ਹਿਸਾਬ}} — \emph{hisāb} — \textquotedblleft{}a reckoning; an account of money\textquotedblright{}.

\begin{quote}
\pa{ਦੁਕਾਨ} \pa{ਦਾ} \pa{ਹਿਸਾਬ} \pa{ਸਹੀ} \pa{ਹੈ।}
\end{quote}

\begin{quote}
\emph{dukān dā hisāb sahī hai.}
\end{quote}

\begin{quote}
The shop's accounts are correct.
\end{quote}

A masculine noun.

\begin{cousinweb}
\textbf{\pa{ਹਿਸਾਬ}} (\emph{hisāb}) is a loan, usually traced through Persian to Arabic \emph{hisāb}, \textquotedblleft{}counting, reckoning\textquotedblright{}.
\end{cousinweb}

\begin{grammarlens}[title={What you've built}]
One more word for this chapter.
\end{grammarlens}

\subsection*{Your turn}
\begin{itemize}
\raggedright
  \item \emph{Say it:} \emph{hisāb}
  \item \emph{Say it:} \emph{hisāb}, once more
  \item \emph{Say it:} \emph{dukān dā hisāb sahī hai.} — The shop's accounts are correct.
  \item \emph{Recall:} say the Punjabi for a thing, then the Punjabi for goods, then say \emph{hisāb} again
\end{itemize}

\begin{tcolorbox}[breakable,colback=teal!4,colframe=teal!35!black,title={Before you move on}]
What does \pa{ਹਿਸਾਬ} mean? (\textquotedblleft{}A reckoning; an account of money\textquotedblright{}.) Say it once more.
\end{tcolorbox}
\section[\texorpdfstring{\pa{ਲੱਦਣਾ} (laddṇā)}{ਲੱਦਣਾ (laddṇā)}]{\pa{ਲੱਦਣਾ} (laddṇā) — to load}

\label{lesson:PA-C278-laddna}



\begin{quote}
Before the new one: say the Punjabi for goods, then the Punjabi for a reckoning.
\end{quote}

\subsection*{You'll want to know: \pa{ਲੱਦਣਾ}}
\textbf{\pa{ਲੱਦਣਾ}} — \emph{laddṇā} — \textquotedblleft{}to load\textquotedblright{}.

\begin{quote}
\pa{ਮੁੰਡਾ} \pa{ਸਾਮਾਨ} \pa{ਬੱਸ} \pa{ਉੱਤੇ} \pa{ਲੱਦਦਾ} \pa{ਹੈ।}
\end{quote}

\begin{quote}
\emph{muṇḍā sāmān bass utte ladddā hai.}
\end{quote}

\begin{quote}
The boy loads the luggage onto the bus.
\end{quote}

Past (I …): \pa{ਮੈਂ} \pa{ਸਾਮਾਨ} \pa{ਲੱਦਿਆ}, \emph{maiṁ sāmān laddiā}.

\begin{cousinweb}
\textbf{\pa{ਲੱਦਣਾ}} (\emph{laddṇā}) is inherited from Sanskrit \emph{lardayati}, \textquotedblleft{}loads\textquotedblright{}; CDIAL lists Punjabi \emph{laddṇā}, \textquotedblleft{}to load\textquotedblright{}, under it (CDIAL 10966).
\end{cousinweb}

\begin{grammarlens}[title={What you've built}]
One more word for this chapter.
\end{grammarlens}

\subsection*{Your turn}
\begin{itemize}
\raggedright
  \item \emph{Say it:} \emph{laddṇā}
  \item \emph{Say it:} \emph{laddṇā}, once more
  \item \emph{Say it:} \emph{muṇḍā sāmān bass utte ladddā hai.} — The boy loads the luggage onto the bus.
  \item \emph{Recall:} say the Punjabi for goods, then the Punjabi for a reckoning, then say \emph{laddṇā} again
  \item \emph{Write it:} copy \textbf{\pa{ਲੱਦਣਾ}} once, checking the whole word
\end{itemize}

\begin{tcolorbox}[breakable,colback=teal!4,colframe=teal!35!black,title={Before you move on}]
What does \pa{ਲੱਦਣਾ} mean? (\textquotedblleft{}To load\textquotedblright{}.) Say it once more.
\end{tcolorbox}
\section[\texorpdfstring{\pa{ਵਿਕਣਾ} (vikṇā)}{ਵਿਕਣਾ (vikṇā)}]{\pa{ਵਿਕਣਾ} (vikṇā) — to be sold; to sell}

\label{lesson:PA-C278-vikna}



\begin{quote}
Before the new one: say the Punjabi for a reckoning, then the Punjabi for to load.
\end{quote}

\subsection*{You'll want to know: \pa{ਵਿਕਣਾ}}
\textbf{\pa{ਵਿਕਣਾ}} — \emph{vikṇā} — \textquotedblleft{}to be sold; to sell\textquotedblright{}.

\begin{quote}
\pa{ਬਾਜ਼ਾਰ} \pa{ਵਿੱਚ} \pa{ਸਸਤੇ} \pa{ਕੱਪੜੇ} \pa{ਵਿਕਦੇ} \pa{ਹਨ।}
\end{quote}

\begin{quote}
\emph{bāzār vich saste kappṛe vikde han.}
\end{quote}

\begin{quote}
Cheap clothes sell in the market.
\end{quote}

Past (it …): \pa{ਉਹ} \pa{ਪੁਰਾਣਾ} \pa{ਘਰ} \pa{ਵਿਕਿਆ}, \emph{oh purāṇā ghar vikiā}.

\begin{cousinweb}
\textbf{\pa{ਵਿਕਣਾ}} (\emph{vikṇā}) is inherited from Sanskrit \emph{vikrīyate}, \textquotedblleft{}is sold\textquotedblright{}; CDIAL lists Punjabi \emph{vikṇā}, \textquotedblleft{}to be sold\textquotedblright{}, under it (CDIAL 11642).
\end{cousinweb}

\begin{grammarlens}[title={What you've built}]
One more word for this chapter.
\end{grammarlens}

\subsection*{Your turn}
\begin{itemize}
\raggedright
  \item \emph{Say it:} \emph{vikṇā}
  \item \emph{Say it:} \emph{vikṇā}, once more
  \item \emph{Say it:} \emph{bāzār vich saste kappṛe vikde han.} — Cheap clothes sell in the market.
  \item \emph{Recall:} say the Punjabi for a reckoning, then the Punjabi for to load, then say \emph{vikṇā} again
  \item \emph{Write it:} copy \textbf{\pa{ਵਿਕਣਾ}} once, checking the whole word
\end{itemize}

\begin{tcolorbox}[breakable,colback=teal!4,colframe=teal!35!black,title={Before you move on}]
What does \pa{ਵਿਕਣਾ} mean? (\textquotedblleft{}To be sold; to sell\textquotedblright{}.) Say it once more.
\end{tcolorbox}
\section[\texorpdfstring{\pa{ਪੈਕਟ} (paikaṭ)}{ਪੈਕਟ (paikaṭ)}]{\pa{ਪੈਕਟ} (paikaṭ) — a packet}

\label{lesson:PA-C278-paikat}



\begin{quote}
Before the new one: say the Punjabi for to load, then the Punjabi for to be sold.
\end{quote}

\subsection*{You'll want to know: \pa{ਪੈਕਟ}}
\textbf{\pa{ਪੈਕਟ}} — \emph{paikaṭ} — \textquotedblleft{}a packet\textquotedblright{}.

\begin{quote}
\pa{ਮੈਂ} \pa{ਚਾਹ} \pa{ਦਾ} \pa{ਇੱਕ} \pa{ਪੈਕਟ} \pa{ਖ਼ਰੀਦਾਂਗਾ।}
\end{quote}

\begin{quote}
\emph{maiṁ chāh dā ikk paikaṭ kharīdāṅgā.}
\end{quote}

\begin{quote}
I will buy a packet of tea.
\end{quote}

A masculine noun.

\begin{cousinweb}
\textbf{\pa{ਪੈਕਟ}} (\emph{paikaṭ}) is borrowed from English \emph{packet}.
\end{cousinweb}

\begin{grammarlens}[title={What you've built}]
One more word for this chapter.
\end{grammarlens}

\subsection*{Your turn}
\begin{itemize}
\raggedright
  \item \emph{Say it:} \emph{paikaṭ}
  \item \emph{Say it:} \emph{paikaṭ}, once more
  \item \emph{Say it:} \emph{maiṁ chāh dā ikk paikaṭ kharīdāṅgā.} — I will buy a packet of tea.
  \item \emph{Recall:} say the Punjabi for to load, then the Punjabi for to be sold, then say \emph{paikaṭ} again
  \item \emph{Write it:} copy \textbf{\pa{ਪੈਕਟ}} once, checking the whole word
\end{itemize}

\begin{tcolorbox}[breakable,colback=teal!4,colframe=teal!35!black,title={Before you move on}]
What does \pa{ਪੈਕਟ} mean? (\textquotedblleft{}A packet\textquotedblright{}.) Say it once more.
\end{tcolorbox}
\section[\texorpdfstring{\pa{ਜੋੜਨਾ} (joṛnā)}{ਜੋੜਨਾ (joṛnā)}]{\pa{ਜੋੜਨਾ} (joṛnā) — to add up; to join}

\label{lesson:PA-C278-jorna}



\begin{quote}
Before the new one: say the Punjabi for to be sold, then the Punjabi for a packet.
\end{quote}

\subsection*{You'll want to know: \pa{ਜੋੜਨਾ}}
\textbf{\pa{ਜੋੜਨਾ}} — \emph{joṛnā} — \textquotedblleft{}to add up; to join\textquotedblright{}.

\begin{quote}
\pa{ਮੈਂ} \pa{ਕਾਗ਼ਜ਼} \pa{ਉੱਤੇ} \pa{ਸਭ} \pa{ਪੈਸੇ} \pa{ਜੋੜਦਾ} \pa{ਹਾਂ।}
\end{quote}

\begin{quote}
\emph{maiṁ kāghaz utte sabh paise joṛdā hāṁ.}
\end{quote}

\begin{quote}
I add up all the money on paper.
\end{quote}

Past (I …): \pa{ਮੈਂ} \pa{ਹਿਸਾਬ} \pa{ਜੋੜਿਆ}, \emph{maiṁ hisāb joṛiā}.

\begin{cousinweb}
\textbf{\pa{ਜੋੜਨਾ}} (\emph{joṛnā}) is inherited from a reconstructed \emph{*yuṭati}, \textquotedblleft{}is joined\textquotedblright{}; CDIAL lists Punjabi \emph{joṛṇā}, \textquotedblleft{}to join\textquotedblright{}, under it (CDIAL 10496).
\end{cousinweb}

\begin{grammarlens}[title={What you've built}]
One more word for this chapter.
\end{grammarlens}

\subsection*{Your turn}
\begin{itemize}
\raggedright
  \item \emph{Say it:} \emph{joṛnā}
  \item \emph{Say it:} \emph{joṛnā}, once more
  \item \emph{Say it:} \emph{maiṁ kāghaz utte sabh paise joṛdā hāṁ.} — I add up all the money on paper.
  \item \emph{Recall:} say the Punjabi for to be sold, then the Punjabi for a packet, then say \emph{joṛnā} again
  \item \emph{Write it:} copy \textbf{\pa{ਜੋੜਨਾ}} once, checking the whole word
\end{itemize}

\begin{tcolorbox}[breakable,colback=teal!4,colframe=teal!35!black,title={Before you move on}]
What does \pa{ਜੋੜਨਾ} mean? (\textquotedblleft{}To add up; to join\textquotedblright{}.) Say it once more.
\end{tcolorbox}
